\documentclass[11pt]{article}

\usepackage[T1]{fontenc}
\usepackage[utf8]{inputenc}
\usepackage{lmodern}
\usepackage[a4paper,margin=28mm]{geometry}
\usepackage{microtype}
\usepackage{amsmath,amssymb,amsthm,mathtools}
\usepackage{enumitem}
\usepackage{booktabs,tabularx}
\usepackage[colorlinks=true,linkcolor=blue,citecolor=blue,urlcolor=blue]{hyperref}

\newtheorem{theorem}{Theorem}[section]
\newtheorem{proposition}[theorem]{Proposition}
\newtheorem{corollary}[theorem]{Corollary}
\theoremstyle{definition}
\newtheorem{definition}[theorem]{Definition}
\newtheorem{assumption}[theorem]{Assumption}
\newtheorem{example}[theorem]{Example}
\theoremstyle{remark}
\newtheorem{remark}[theorem]{Remark}

\newcommand{\card}{\operatorname{card}}
\newcommand{\vol}{\operatorname{vol}}
\newcommand{\tr}{\operatorname{tr}}

\title{Coherence-Selection Events and Causal Sets in Modal Triplet Theory:\\
\large Non-Zeno Conditions, Reset Laws, and Covariant Statistics}
\author{Peter Nero}
\date{July 2026\\Version 2}

\begin{document}
\maketitle

\begin{abstract}
Coherence-selection dynamics can generate an effective causal set only when
the event law prevents accumulation and is covariantly defined. This paper
states two sufficient routes. In a deterministic hybrid model, a reset
margin together with a uniform bound on trigger growth gives a positive
dwell time and hence finitely many events on every bounded trajectory
interval. In a stochastic model, finite compensator measure on every compact
spacetime region implies almost-sure local finiteness. Once local finiteness
is established, events inherit a partial order from a globally hyperbolic
spacetime. Positive stability margins alone do not provide either result,
and noninvertible projection alone does not define event times, reset states,
outcome probabilities, or irreversibility. Lorentz-invariant statistics
require a concrete covariant event law, such as a scalar intensity with
respect to spacetime volume; they do not follow from coherence language.
The current MTT corpus does not yet select the required trigger, reset, or
hazard. The construction is therefore a rigorous conditional event
encoding, not a derivation of fundamental spacetime discreteness or
measurement outcomes.
\end{abstract}

\section*{Version 2 Revision Note}
\begin{description}[leftmargin=3.2cm,style=nextline]
\item[Supersedes] Version 1.0, \emph{Causal Sets as Event-Selection Shadows
of Coherence Breakdown}.
\item[Reason] The original paper inferred local finiteness from positive
stability margins, left the transition/reset law undefined, and claimed
approximately Lorentz-invariant statistics without a stochastic source.
\item[Resolution] Deterministic dwell-time and stochastic compensator
criteria are proved, the event/reset data are typed explicitly, and
covariance is made a property of the selected law.
\item[Retained result] Dynamically generated, locally finite events inherit
the causal order of the effective spacetime and can be represented as a
causal set.
\item[Remaining boundary] MTT must select the trigger or hazard, reset law,
spacetime density, and any outcome instrument from the same physical branch.
\end{description}

\tableofcontents

\section{Events need a law}

An ``event'' is not created by naming a loss of coherence. A dynamical event
model must specify:
\begin{itemize}
\item a state space and continuous evolution between events;
\item an event trigger or stochastic hazard;
\item the spacetime location assigned to an event;
\item a reset or transition rule after the event;
\item whether the event is an observed record, an objective state change, or
only a coarse-graining marker.
\end{itemize}

These distinctions matter in MTT. A noninvertible lower description can
discard information without selecting a unique physical outcome. Likewise,
a stable fixed point can exist without any finite sequence of transition
times.

\section{Deterministic hybrid route}

Let \(z(t)\) be a continuous state between event times
\[
 0<t_1<t_2<\cdots.
\]
Let \(h(z)\) be a scalar trigger. An event occurs when \(h\) reaches zero
from below, after which a reset map \(R\) replaces \(z(t_n^-)\) by
\(z(t_n^+)=R(z(t_n^-))\).

\begin{theorem}[Uniform dwell-time criterion]
\label{thm:dwell}
Assume that after every event
\[
 h(z(t_n^+))\leq-\delta
\]
for a fixed \(\delta>0\), and that between events
\[
 \left|\frac{d}{dt}h(z(t))\right|\leq M
\]
for a fixed finite \(M>0\). Then
\[
 t_{n+1}-t_n\geq\frac{\delta}{M}.
\]
Consequently a trajectory has at most
\[
 1+\frac{MT}{\delta}
\]
events in any interval of duration \(T\).
\end{theorem}

\begin{proof}
After the reset, the trigger must increase by at least \(\delta\) before it
can reach zero. The derivative bound implies that this takes at least
\(\delta/M\). Summing the minimum separations gives the counting bound.
\end{proof}

This is the missing non-Zeno theorem. A positive stability margin for a
fixed point does not imply either a reset margin or a bound on trigger
growth.

\begin{example}[Stability without a dwell-time theorem]
The continuous system \(\dot x=-x\) has a positive linear stability rate at
the origin. If one separately declares events at \(t_n=1-2^{-n}\), the
events accumulate at \(t=1\). The flow's stability has not prevented the
Zeno schedule because the event law was independent of the flow.
\end{example}

\section{Stochastic hazard route}

Let \(N(K)\) count events in a measurable spacetime region \(K\). A point
process may be specified through a compensator or intensity measure
\(\Lambda\).

\begin{theorem}[Finite-compensator local finiteness]
\label{thm:compensator}
Suppose an event point process satisfies
\[
 \mathbb E[N(K)]=\Lambda(K)<\infty
\]
for every compact \(K\subset Y_4\). Then \(N(K)<\infty\) almost surely for
every member of a countable compact exhaustion. In particular, the event set
is locally finite almost surely.
\end{theorem}

\begin{proof}
A nonnegative extended-integer random variable that equals infinity with
positive probability has infinite expectation. Therefore finite expectation
implies \(N(K)<\infty\) almost surely. Apply this to a countable compact
exhaustion.
\end{proof}

A sufficient form is
\[
 \Lambda(K)=\int_K\lambda(x,z)\,d\vol_g(x),
\]
with a nonnegative scalar predictable intensity \(\lambda\) that is locally
integrable. A uniform compact-region bound on \(\lambda\) gives finite
\(\Lambda(K)\).

This theorem does not derive Poisson statistics. Correlated point processes,
self-exciting processes, and state-dependent hazards can all satisfy local
finiteness. Their physical predictions differ.

\section{From events to a causal set}

Let \(C\subset Y_4\) be the generated event set. On a globally hyperbolic
spacetime define
\[
 x\preceq y
 \quad\Longleftrightarrow\quad
 y\in J^+(x).
\]
The kinematic companion paper proves that any locally finite \(C\) with
this inherited order is a causal set. Thus either
\autoref{thm:dwell}, with the required spatial population bound, or
\autoref{thm:compensator} can provide the local-finiteness input.

For a deterministic family of trajectories, a compact spacetime region must
intersect only finitely many relevant trajectories, or a separate uniform
spatial density bound is needed. Dwell time along each of infinitely many
worldlines would not by itself control the total number of events in the
region.

\section{Reset, irreversibility, and records}

A reset map can be invertible or noninvertible. If it is noninvertible, the
effective hybrid dynamics loses information. That is a model of
irreversibility at the chosen descriptive level, not automatically a
fundamental arrow of time.

For quantum states, an outcome event requires an instrument
\(\{\mathcal I_a\}\), not only an effect or decoherence channel. The
probability and conditional state are
\[
 p_a=\tr[\mathcal I_a(\rho)],
 \qquad
 \rho_a=\frac{\mathcal I_a(\rho)}{p_a}
\]
when \(p_a>0\). A nonselective channel
\(\sum_a\mathcal I_a\) does not say which event occurred. The current MTT
Born-source result is closed only on a restricted selected recorder and
remains open for arbitrary apparatus contexts and objective ontic
selection.

Measurement has no privileged metaphysical role here. It is one physical
process capable of producing records. The same event mathematics can
describe switching, threshold crossings, detector clicks, or other
transitions.

\section{Covariance and Lorentz statistics}

Lorentz-invariant or generally covariant statistics are properties of the
event law. In Minkowski spacetime, a constant scalar Poisson intensity
\[
 \Lambda(B)=\rho\int_Bd^4x
\]
is invariant in law under proper orthochronous Poincare transformations. In
curved spacetime, a scalar \(\lambda(x,z)\) integrated against
\(d\vol_g\) is covariantly defined if the state variables and evolution are
also geometrically specified.

By contrast, a reset at equal coordinate-time intervals or a hazard tied to
a preferred foliation generally breaks Lorentz symmetry. An approximately
isotropic numerical event cloud is not a proof of invariant law.

\section{Relation to MTT data}

The local \(1+2+3\) filtration, shared phase line, finite q79 Hessian, and
fixed-point stability results do not currently emit:
\begin{itemize}
\item a four-dimensional scalar trigger \(h\);
\item a reset map \(R\);
\item a spacetime hazard \(\lambda\);
\item a covariant event-location rule;
\item an event density in physical units.
\end{itemize}

These objects could be derived in a future same-source event model. For
example, a response functional might define a scalar threshold, while a
selected detector or basin map defines the reset. But the cross-sector
intertwiner and unit map must be exhibited. An internal spectral gap alone
does not supply a four-volume intensity.

\section{Status ledger}

\begin{center}
\begin{tabularx}{\textwidth}{@{}
>{\raggedright\arraybackslash}p{0.31\textwidth}
>{\raggedright\arraybackslash}p{0.18\textwidth}
>{\raggedright\arraybackslash}X@{}}
\toprule
Statement & Status & Boundary\\
\midrule
Reset margin plus derivative bound gives dwell time & Exact &
\autoref{thm:dwell}.\\
Finite compensator gives local finiteness & Exact &
\autoref{thm:compensator}.\\
Local events inherit causal-set order & Exact, imported &
Kinematic companion theorem.\\
Positive stability margin excludes Zeno events & False in general &
Event/reset law is additional.\\
MTT selects trigger/reset law & Open &
No same-source event generator.\\
MTT predicts Lorentz-invariant event statistics & Open &
Requires a covariant stochastic or hybrid law.\\
Event channel selects objective outcome & Open &
Instrument, probability, and ontology are separate.\\
\bottomrule
\end{tabularx}
\end{center}

\section{Discussion}

The corrected event route is sharper than the original proposal. It gives
two concrete mathematical exits from the local-finiteness blocker and makes
clear what a simulation would need to test. A hybrid simulation should log
the reset margin and maximum trigger derivative. A stochastic simulation
should estimate or certify compact-region compensators and test covariance
of the full law, not only the appearance of one sample.

The shortest path to an MTT-specific result is to construct one physical
event source on the already selected free q79 QFT net. A detector
instrument, a covariant hazard derived from its response kernel, and a
finite-compensator proof would produce a genuine event causal set without
assuming Poisson sprinkling.

\section{Conclusion}

Coherence-selection events can form a causal set, but only after
non-accumulation and causal-order conditions are proved. Uniform reset
margins and trigger bounds provide one route; finite stochastic compensators
provide another.

MTT has not yet selected either event law. The present result is therefore a
conditional mathematical framework. It replaces an unsupported local-
finiteness claim with exact, testable criteria and preserves the distinction
between effective irreversibility, observed records, and objective outcomes.

\begin{thebibliography}{99}

\bibitem{Bombelli}
L.~Bombelli, J.~Lee, D.~Meyer, and R.~D.~Sorkin,
``Space-time as a causal set,''
\emph{Physical Review Letters} \textbf{59} (1987), 521--524.

\bibitem{DaleyVereJones}
D.~J.~Daley and D.~Vere-Jones,
\emph{An Introduction to the Theory of Point Processes, Volume I},
Springer, 2003.

\bibitem{Goebel}
R.~Goebel, R.~G.~Sanfelice, and A.~R.~Teel,
\emph{Hybrid Dynamical Systems: Modeling, Stability, and Robustness},
Princeton University Press, 2012.

\bibitem{MTTCausalKinematic}
P.~Nero,
\emph{Causal Sets as a Conditional Coarse-Graining of Modal Triplet
Theory},
current revised MTT paper corpus, 2026.

\bibitem{MTTMeasurement}
P.~Nero,
\emph{Measurement and Coherence Selection in Modal Triplet Theory},
current revised MTT paper corpus, 2026.

\end{thebibliography}

\end{document}
